\section{アプローチ07（重複排除手法）}

\subsection{手法の概要}
アプローチ06で顕在化した「同一人物に対する複数の重複したバウンディングボックスの生成および過剰な計算資源の消費」という課題を解決するため、アプローチ07ではポーズ検出結果の骨格ポイントを活用した再切り出しと、領域の重なりに基づく重複排除（Deduplication）手法を導入した。

検出結果の正常・不正常に関わらず、ポーズ検出によって抽出された関節キーポイント群に対して一定のマージンを付与した新たなバウンディングボックスを再定義する。これにより、元領域内で最も目立った人物のみを正確な矩形として抽出するとともに、余白が大きく不安定であった検出領域の補正を行った。

最後に、残された各人物バウンディングボックス間の面積重複度（Intersection over Union: IoU）を計算し、重複が確認された場合には「別人物」ではなく「同一人物のより描画状態が良い（写りが良い）大きい方のバウンディングボックス」を保持する判定ロジックを適用した。

\subsection{結果およびアプローチ06との比較・課題}

\subsubsection{アプローチ06との比較}
骨格情報を基準としたバウンディングボックスの再構成とIoUに基づく重複判定の組み合わせにより、単一の人物に対して乱立していた不要な検出枠が大幅に整理された。また、余白が過剰に含まれていた低精度なバウンディングボックスが適正な範囲へ補正され、後続処理へのノイズ伝播を防ぐことが可能となった。

\subsubsection{判明した課題}
\begin{enumerate}
    \item \textbf{垂直（縦）方向における誤断片化・乱立の残存}：
    二次元のIoU重なり判定のみでは、待機列特有の前後重なりにおいて、ポーズ検出を通過した比較的小さなバウンディングボックスが縦方向に連続して乱立し、すり抜けて残ってしまう現象が確認された。
    \item \textbf{エッジ環境における推論処理のメモリ蓄積}：
    重複排除の前段として多数の推論処理を反復実行するため、Webブラウザ環境（特にWebGLレンダリングバックエンド）においてテンソルやテクスチャのキャッシュが蓄積し、持続的な実行時にメモリ消費が増大するリスクが判明した。
\end{enumerate}
